\documentclass[11pt]{article}
\usepackage[margin=1.1in]{geometry}
\usepackage{amsmath,amssymb,amsthm}
\usepackage{booktabs}
\usepackage{array}
\usepackage[hidelinks]{hyperref}
\usepackage{xcolor}
\usepackage{parskip}

\newcommand{\CNOT}{\textsc{cnot}}
\newcommand{\code}[1]{\texttt{#1}}
\newcommand{\gf}{\mathrm{GF}(2)}
\newcommand{\pe}{\mathrel{{+}{=}}}    % t += e  means  t <- t XOR e over GF(2)
\newtheorem{lemma}{Lemma}
\newtheorem{definition}{Definition}

\title{\textbf{A Nonlinear Share Encoding by Deferred Mask Compensation}\\[2pt]
\large Design note: the split-RG3 move (RG4), mask towers, and the
maintenance discipline that removes the degree-1 snapshot leak}
\author{local\_mixing project}
\date{July 21, 2026 \\ \small branch \code{ssg-gen-mix-clean} ---
implemented same day (\code{MaskLedger} in \code{src/replace/gadgets.rs},
flags on \code{sss --cnot}); masks default OFF until the \S10 grid
validates}

\begin{document}
\maketitle

\begin{abstract}
\noindent The two-share gadget carries every logical value $v$ as an XOR
pair $v = c_0 \oplus c_1$. This makes $v$ a \emph{degree-1 function of the
physical wires at every snapshot} of the circuit, and the measured
consequence is the persistent computational-progress diagonal in the
\code{hmap\_affine} reconstruction maps: it survives 4M fmix moves with
alignment intact ($\rho = 1.00$, perm $z \approx 9.9$), mixing only
shallows it. Nonlinear RGs are provably orthogonal to this leak --- it is
the \emph{encoding}, not the transitions. This note designs the fix inside
the existing vocabulary and wire budget: a \emph{deferred-mask encoding} in
which quadratic mask terms are XORed into one carrier of a pair and
compensated only later, so that between injection and compensation the
value reads as $v = c_0 \oplus c_1 \oplus u{\wedge}w$ --- degree 2 in the
wires, with random support. The move is exactly the existing RG3
compensated pair-write \emph{split in time} (hence ``split-RG3''; we call
it RG4). Stacked masks drive the best affine readout to $1/2$ by
piling-up; mask \emph{towers} (masks whose sources are themselves
overwritten by registered quadratic writes) push snapshot readout to
degree 3 and beyond. We give the invariant, the gate-level realizations
(polymorphic, census-neutral, no bare X), the maintenance discipline, the
security arithmetic, a cost model, and the verification and measurement
plan.
\end{abstract}

\section*{Implementation status (v1, 2026-07-21)}

Implemented in \code{src/replace/gadgets.rs} (\code{MaskLedger},
\code{WirePoly}, \code{emit\_poly\_add}) on both \code{sss --cnot}
gadgetizer paths, flags \code{--mask-cov/--mask-k/--mask-depth/--mask-taper}
(default off). The move is realized as a \emph{split nonlinear RG3 over
logical source values}: a mask on value $v$ XORs the RG3 randomizer
$v_{r1} \vee \lnot v_{r2}$ (two \emph{value} sources --- degree~2, the
quadratic $v_{r1} v_{r2}$ being what makes $v$ degree-2) into one carrier
at inject and re-applies it at flush. Sourcing on logical \emph{values}
(rebuilt from their current carriers each emit) makes the mask invariant
under all RG churn --- RGs preserve logical values --- so the only
disturbance is a CG that recomputes a source, whose schedule is the input
gate list's target column, known in advance. \textbf{Lookahead} picks the
two unmasked values with the most distant next recomputation, so inject and
flush are widely separated instead of collapsing back into a plain RG3
(wire-sourcing did collapse: at $n{=}8$ it took 139 injects and peaked at
0.67 coverage, vs.\ 54 and 1.00 for value-sourcing). When the virtual
circuit \emph{uses} $v$ as a control, the mask is \textbf{peeked}: un-mask,
run the vanilla CG on the true value, re-mask --- the mask keeps its
identity and lifetime (the simple form of the mask-aware CG,
\S\ref{sec:disc}), at the cost of a one-step exposure of $v$ where the
original already uses it. The \emph{cascade-free invariant} (now at value
level: masks sourced only on unmasked values) is retained. Two things stay
deferred: in-flight towers (\S\ref{sec:towers}, masks stay degree~2) and
$\text{coverage}{=}1$ (a source pool must stay unmasked;
\code{--mask-cov}${\sim}0.75$). Correctness: exhaustive
\code{masked\_gadgetize\_*}, \code{mask\_peek\_bracket\_*}, and
\code{mask\_inject\_then\_flush\_*} tests plus the driver's pinned
zero-slice check.

\section{The problem: the snapshot leak}

The degree-bounded reconstruction adversary (\code{hmap\_affine}) fits
each bit of the original circuit's state after prefix $i$ as a degree-$d$
$\gf$ function of \emph{all} wires of the gadgetized/mixed circuit after
prefix $j$. Under the XOR 2-sharing every logical value is
$v = c_0 \oplus c_1$: a degree-1 function of the snapshot, at every $j$,
no matter how nonlinear the \emph{transitions} between snapshots are. The
ridge-measure campaign made the consequence quantitative: a progress
diagonal with $\rho = 1.00$ and perm $z \approx 9.9$ survives 4M fmix
moves; mixing only shallows its depth ($0.35 \to 0.27$). The 2026-07-20b
nonlinear-RG revision explicitly documented this residue: nonlinear RGs
make \emph{crossing} the body expensive for low-degree predictors, but a
snapshot \emph{readout} needs to cross nothing. The open-thread list names
the two real fixes: a nonlinear share encoding, and heavier global
reordering. This note is the first.

Scope: this design attacks the \emph{snapshot readout} leak only. Gate
\emph{order} legibility (positional structure) remains the job of the
commuting shuffle and the mixing phase; the static/syntactic census front
is a separate instrument. Both are noted where they interact.

\section{Why the pair itself cannot be nonlinear}

\begin{lemma}
Every balanced function $D : \{0,1\}^2 \to \{0,1\}$ is affine.
\end{lemma}
\begin{proof}
Balanced means $|D^{-1}(1)| = 2$. There are $\binom{4}{2} = 6$ such
functions: $c_0$, $\bar c_0$, $c_1$, $\bar c_1$, $c_0 \oplus c_1$,
$\overline{c_0 \oplus c_1}$ --- all affine.
\end{proof}

A per-pair decode map must be balanced, or the carriers are individually
biased toward $v$ and leak by correlation alone. So with two carriers per
value there is \emph{no} nonlinear uniform encoding: nonlinearity must
come from (i) more carriers per value, (ii) accepted bias, or (iii)
\emph{cross-pair state}. Option (i) breaks the $\times 2$ wire contract
(all $2n$ wires are already pair carriers); option (ii) is a leak by
construction. This design is option (iii): the decode of value $i$
references wires outside pair $i$, transiently and by registered
bookkeeping.

\section{The deferred-mask invariant}

\begin{definition}[Masked sharing]
The gadget state carries, per logical value $i$, its carrier pair
$(c_{i,0}, c_{i,1})$ \emph{and} a (possibly empty) list of pending mask
polynomials $\varphi_{i,1}, \dots, \varphi_{i,k}$, each a small
multilinear polynomial over the \emph{current} contents of a few
registered source wires. The invariant, restored at every gadget-unit
boundary (see \S5), is
\[
v_i \;=\; c_{i,0} \oplus c_{i,1} \oplus \bigoplus_t \varphi_{i,t}
(\text{current wires}).
\]
\end{definition}

\textbf{RG4 is RG3 split in time.} The existing RG3 writes
$\varphi = r_1 \vee \bar r_2$ into \emph{both} carriers of one pair,
back-to-back: the value is untouched at every subsequent snapshot, which
is why RG3 re-randomizes carriers but cannot mask the value. RG4 emits the
same two writes \emph{separated by an arbitrary stretch of body}:

\begin{itemize}
\item \textbf{RG4-inject}: choose value $i$ (coverage census, \S9),
a carrier of pair $i$, source wires $u, w \notin \{c_{i,0}, c_{i,1}\}$,
and a realization (\S4) of a polynomial $\varphi$ with quadratic content
$u w$. Emit it targeting the chosen carrier; register $\varphi$ on value
$i$.
\item \textbf{RG4-flush}: emit a realization of the \emph{same} $\varphi$
targeting \emph{either} carrier of pair $i$'s \emph{current} pair, then
unregister. (Both carriers XOR into $v_i$, so compensation on the other
carrier is equally correct; the junk output $J$ differs, which is free.)
\item \textbf{RG4-extend} (\S6): when a registered write is about to hit a
source of $\varphi$, optionally rewrite the pending polynomial in terms of
the post-write wires instead of flushing --- this is what builds towers.
\end{itemize}

Between inject and flush, $v_i$ is a degree-2 function of the wires whose
quadratic monomial $u w$ has \emph{random support} among all $2n$ wires.

\textbf{Bookkeeping under RG1/RG2.} Masks are registered by \emph{value
index}. The RG1/RG2 networks are verified (existing exhaustive tests) to
transport pair-XOR values for \emph{all} wire contents, and the emitters'
\code{pairs} updates keep \code{pairs[$i$]} = the carriers currently
holding value $i$'s (masked) content. So a mask attached to value $i$
follows it automatically through value-swaps and re-pairings; a flush
simply targets a carrier of the \emph{current} \code{pairs[$i$]}. Sources,
by contrast, are registered as \emph{physical wires}: the interception
discipline (\S5) is defined at wire level and needs no relocation logic.

\section{Gate-level realizations: polymorphic, no bare X, census-neutral}

Vocabulary: g57s, conjunction fragments with one or two controls of any
polarity, positive and negative \CNOT{}s; never a bare X. The primitive
``adders'' available for a target $t$ and sources $u, w$, as polynomials
over $\gf$:

\begin{center}
\begin{tabular}{llll}
\toprule
gate & adds to $t$ & degree & note \\
\midrule
$t \pe u$ (\CNOT) & $u$ & 1 & \\
$t \pe \bar u$ ($\lnot$\CNOT) & $1 \oplus u$ & 1 & folds a constant \\
$t \pe (u^{\varepsilon} w^{\delta})$ (fragment) & product of literals & 2 & \\
$t \pe (u \vee \bar w)$ (g57) & $1 \oplus w \oplus u w$ & 2 & the RG3 shape \\
\bottomrule
\end{tabular}
\end{center}

Any pending polynomial with monomials in $\{1, u, w, uw\}$ is a $\gf$-sum
of these. The ledger records $\varphi$ \emph{as a polynomial}; inject and
flush each draw an \emph{independent} realization. For the pure monomial
$\varphi = uw$, six of them:

\begin{center}
\small
\begin{tabular}{clc}
\toprule
\# & realization (order free) & gates \\
\midrule
1 & $t \pe (u\,w)$ & 1 \\
2 & $t \pe (u \vee \bar w)$, \quad $t \pe \bar w$ & 2 \\
3 & $t \pe (w \vee \bar u)$, \quad $t \pe \bar u$ & 2 \\
4 & $t \pe (\bar u\, w)$, \quad $t \pe w$ & 2 \\
5 & $t \pe (u\, \bar w)$, \quad $t \pe u$ & 2 \\
6 & $t \pe (\bar u\, \bar w)$, $t \pe \bar u$, $t \pe w$ & 3 \\
\bottomrule
\end{tabular}
\end{center}

Three design consequences:
\begin{itemize}
\item \textbf{No pairing signature.} A na\"ive scheme (flush $=$ the
identical gate again) hands the static adversary a free un-masking pass:
scan for equal gates whose pins are un-written between them and cancel.
With independent draws, inject and flush share no shape, no polarity
pattern, and possibly not even a target wire (either-carrier flushing),
and same-shape re-occurrence is common in an ordinary g57-rich body.
\item \textbf{No bare X.} Constants fold into $\lnot$\CNOT{}s and
complemented literals, exactly as in the CG menu.
\item \textbf{Census neutrality and DB-warmth.} The material is
RG3-like: g57s with sprinkled $\pm$\CNOT{}s and fragments --- the same
mix the CG menu already emits, and g57-dominant, hence warm for the
frozen-store phase-A mixing. Realization draws can be biased toward the
g57-bearing variants (2, 3) to keep the body's g57 fraction high.
\end{itemize}

Sources may be any wires outside pair $i$ (all wires are carriers of some
pair; that is fine and is what makes the mask supports random). Preferring
sources on \emph{distinct} pairs avoids concentrating probe information;
preferring sources with a distant next write (\S8) makes masks cheap.

\section{The maintenance discipline}\label{sec:disc}

Emission proceeds, as today, in \emph{units}: one CG (menu draw), one RG
(uniform nonlinear draw), one W/bookend gate, one RG4 event. Each unit has
a set of \emph{net writes} --- wires whose content differs after the unit
--- which excludes borrowed-and-restored scratch and the transient
collapse writes inside CG variants 0--2. (Inside a unit the invariant may
be transiently false; it is asserted only at unit boundaries. Deferred
compensation is unaffected: correctness needs source values at flush time
to equal source values at inject time, and restored excursions in between
are invisible.)

\textbf{Interception rule.} Before emitting any unit whose net writes
intersect the source set of a pending $\varphi$ on value $i$: either
\emph{flush} $\varphi$ first, or --- when the write is itself registered
with a known added polynomial $\psi$ (an RG3 write, another RG4 inject, a
linear RG2 \CNOT) and the tower cap allows --- \emph{extend} (\S6).
Un-registered complex writes (a CG landing $f$ on its target carrier)
always force a flush of masks \emph{sourced} on that carrier.

\textbf{Writes to masked pairs are free.} A CG targeting value $i$ adds
$f$ to one carrier; $v_i' = v_i \oplus f$ holds with the same pending
$\varphi$. RG3 pair-writes and RG1/RG2 relocations are equally transparent
(\S3). Masks ride through all of it.

\textbf{Reads of masked pairs.} A CG \emph{reading} value $i$ (as $B$ or
$C$) would compute with $c_0 \oplus c_1 = v_i \oplus \varphi$. Two
policies:

\begin{itemize}
\item \textbf{(a) flush-and-reinject} (default): flush all pending masks
on the operand values, emit the CG, re-inject fresh masks immediately
after. Simple, strictly small units; coverage has a hole a few gates wide
at each read.
\item \textbf{(b) mask-aware CG}: fold the correction into the CG. For a
single mask $\varphi = \mu$ on $B$ (so $B = c_0^B \oplus c_1^B \oplus
\mu$) the CG's error is $f_{\text{true}} \oplus f_{\text{computed}} = \mu
C$. With a collapse-style variant holding $C$ on $c_0$ during its window,
the degree-3 correction $a^* \pe \mu\,c_0$ is realized in-vocabulary by a
borrowed-and-restored scratch $h$ (any wire outside the unit, not a
pending source):
\[
h \pe (u\,w); \qquad a^* \pe (h\,c_0); \qquad h \pe (u\,w);
\qquad a^* \pe (h\,c_0)
\]
--- net effect $(h \oplus \mu)C \oplus hC = \mu C$, four fragments. A
masked $C$ costs the same with $\mu(1 \oplus B)$; both operands masked, or
multi-term masks: flush down to one term first. Policy (b) keeps coverage
continuous at $+4$ fragments per masked read.
\end{itemize}

\textbf{Decode.} All pending masks are flushed before the $W^{-1}$
curtain. To avoid a syntactic flush burst at the seam, coverage
\emph{tapers}: in the last $L$ gaps (L $\approx$ the mean mask lifetime,
\S8) dying masks are not re-injected, so flushes thin out gradually and
the commuting shuffle disperses the remainder. Symmetrically, coverage
ramps during the first gaps as the census injects. The two forced I/O
corners of the heatmap are expected to stay legible; everything between
should not.

\section{Mask towers: past degree 2}\label{sec:towers}

A single mask makes snapshot readout degree 2. The extension move makes it
deeper. Suppose $\varphi = u w$ is pending and a registered write $u \pe
\psi$ arrives (e.g.\ an RG3 half: $\psi = r_1 \vee \bar r_2 = 1 \oplus r_2
\oplus r_1 r_2$). Extending rewrites the pending polynomial over the
post-write wires: with $u_{\text{new}} = u \oplus \psi$,
\[
\varphi \;=\; u w \;=\; (u_{\text{new}} \oplus \psi)\,w
\;=\; u_{\text{new}} w \,\oplus\, w \,\oplus\, r_2 w \,\oplus\, r_1 r_2 w .
\]
The pending polynomial now contains the cubic monomial $r_1 r_2 w$:
snapshot readout of $v_i$ is degree 3 until the flush. The flush of a
cubic monomial is the same scratch pattern as in \S5(b):
$h \pe (r_1 r_2)$; $t \pe (h w)$; $h \pe (r_1 r_2)$; $t \pe (h w)$ ---
four fragments, plus one or two gates for the residual quadratic and
affine terms. Each further registered write to a surviving source deepens
the tower by one degree and lengthens the eventual flush by one scratch
network.

Rules that keep towers sound and bounded:
\begin{itemize}
\item Only \emph{registered} writes extend (RG3 halves, RG4 injects, RG2's
affine carrier moves --- anything whose added polynomial the emitter
knows). CG target writes never extend; they flush masks sourced on the
target.
\item \code{--mask-depth} caps tower degree (default 2, i.e.\ allow one
extension $\to$ snapshot degree 3). Beyond the cap, the interception rule
falls back to flush.
\item Extension is \emph{opportunistic}: the scheduler does not
manufacture writes to sources; it upgrades collisions that happen anyway
--- RG traffic over $2n$ wires supplies them steadily (\S8).
\item The flush of a depth-$d$ tower costs $O(d)$ scratch networks
($\approx 4d$ fragments); the ledger tracks the exact polynomial, so
correctness is bookkeeping, not search.
\end{itemize}

Degree ladder: depth 0 (plain pair) $= $ degree-1 readout; depth 1 $=$
degree 2; depth $d$ $=$ degree $d{+}1$. The full-degree-2 adversary ---
tractable at $n{=}16$, intractable at production width --- is defeated
from depth 2 up; this is exactly the validation lever in \S10.

\section{What it buys, quantitatively}\label{sec:buys}

\textbf{Degree-1 readout.} A best affine fit absorbs $\varphi$'s affine
part and mismatches on its quadratic monomial: error $\Pr[uw = 1] = 1/4$
per mask (heuristically treating sources as balanced). Stacked masks
compound by piling-up: $k$ independent terms give combined bias
$2^{k-1}(1/4)^k = \tfrac12 2^{-k}$, i.e.\ best affine error
\[
\renewcommand{\arraystretch}{1.1}
\begin{array}{c|ccccc}
k & 1 & 2 & 3 & 4 & 5 \\
\hline
\text{error} & 0.250 & 0.375 & 0.4375 & 0.469 & 0.484
\end{array}
\]
against the ideal $0.5$. The measured pre-mix ridge depth is
$\approx 0.35$: even $k{=}1$ more than halves it; $k{=}3$ puts it near the
permutation-test noise floor. Coverage scales linearly: a bit unmasked at
snapshot $j$ contributes its full readability, so the census must keep
coverage near 1 through the body (\S9).

\textbf{Restricted degree-2.} The feasible quadratic adversary regresses
over a wire subset of size $m$ (full degree-2 at production width,
$\binom{2n}{2} \sim 10^5$ regressors, is intractable --- the campaign's
own finding). A mask term is stripped only if \emph{both} sources land in
the subset: probability $\approx (m/2n)^2$ per term, $2.4\%$ at $m{=}40$,
$2n{=}256$ --- and stripping a $k$-stack needs all $k$. Random-support
masking is, in expectation, as opaque to restricted degree-2 as to
degree 1.

\textbf{Full degree-2 and up.} Depth-1 masks are readable by an
unrestricted degree-2 fit (it can include every $uw$). Towers answer:
depth 2 requires degree 3, and so on. The design makes the depth knob pay
for exactly as much adversary as one wants to defeat.

\textbf{Honest caveats.} (i) Sources are circuit values, not independent
uniform bits; the piling-up arithmetic is a heuristic and the measurement
plan, not the algebra, is the arbiter. (ii) A degree-1 fit may exploit
\emph{correlations}: if $u \wedge w$ happens to be affinely computable
from other wires at snapshot $j$ (e.g.\ a wire that currently holds a
related product), the mask is transparent there; random source draws make
this rare but the heatmap will show whatever survives. (iii) Mask reads
give the probing adversary two-carrier information --- already the
accepted trade of the nonlinear-RG revision. (iv) Snapshot readout of
\emph{carriers} (not values) stays degree 1 by definition; it is the
value readout, which is what \code{hmap\_affine} measures against $C$'s
state, that the masks push up the degree ladder.

\section{Cost model}

Let $\bar w \approx 5$ be the mean net writes per SG gap (CG target $+$ RG
traffic), and $2n = 256$. A pending mask with two random sources dies
(source hit) with probability $\approx 2\bar w / 2n \approx 4\%$ per gap:
mean lifetime $\approx 26$ gaps. With full coverage at $k$ terms per value
($K = kn$ pending), expected maintenance per gap:
\[
\underbrace{k n \cdot \frac{2 \bar w}{2n}}_{\text{deaths}\; \approx\, k
\bar w}
\;+\;
\underbrace{2k}_{\text{CG reads (policy a)}}
\quad \text{events} \;\times\; \approx 2.5 \text{ gates/event}
\;\approx\; 17.5k \text{ gates per gap}
\]
against today's $\approx 10$ (CG $5.1$ $+$ RG $4.7$): body $\times 2.7$
at $k{=}1$, na\"ively. Two levers cut it: \emph{lookahead source
selection} --- the gadgetizer knows the entire future unit stream, so
injecting from wires with distant next net-write turns most deaths into
deliberate end-of-life refreshes (floor: the $2k$ read events, $\approx
{+}50\%$ body at $k{=}1$); and policy (b), which replaces
flush$+$reinject at reads with $+4$ in-place fragments and never opens a
coverage hole. Realistic envelope: \textbf{body $\times 1.5$--$2.7$ at
$k{=}1$, $\times 2.5$--$6$ at $k{=}3$}; the $n{=}16$ benchmark grid
(\S10) prices the curve exactly, in the same spirit as the
generation-vs-size benchmark (whose lesson --- the first sweep is the
whole price --- plausibly recurs here as ``coverage is the whole price,
depth is cheap'').

\section{Knobs, meters, and machinery reuse}

The phase-A generation work built exactly the control pattern this needs,
and it transfers verbatim:

\begin{center}
\small
\begin{tabular}{lll}
\toprule
knob / meter & meaning & analogue \\
\midrule
\code{--mask-cov} & target fraction of values masked (implemented default
0 = off until the \S10 grid validates; production target 1.0) &
\code{--gen-target} \\
\code{--mask-k} & target terms per value (default 1) & --- \\
\code{--mask-depth} & tower degree cap (default 2) & --- \\
\code{--mask-read-policy} & \code{flush} $|$ \code{aware} & --- \\
\code{--mask-taper} & gaps of end-taper (default: mean lifetime) & --- \\
\code{mcov=} & coverage at emission time & \code{cov=} \\
\code{mk=}, \code{mdeep=} & mean stack size, tower histogram & --- \\
\code{mpaid=} & mask-gate ledger (gates spent on inject/flush/extend) &
\code{paid=} \\
\bottomrule
\end{tabular}
\end{center}

An unmasked-value census drives injection exactly as the laggard census
drives DB targeting: rebuilt cheaply per gap, biased draws toward
uncovered values.

\section{Integration}

\textbf{Commuting shuffle: safe without modification.} Correctness of a
deferred pair needs only ``inject before any source net-write before
flush'', and every one of those pairs \emph{collides} in the
\code{XGate::collides} sense (a source write vs.\ a gate reading that
source). The shuffle preserves the relative order of colliding pairs, so
every mask window survives reordering, while the inject/flush gates still
migrate freely within their windows --- which is precisely the
dispersion we want.

\textbf{fmix.} Mask gates are ordinary vocabulary; no bookkeeping
survives into the mixed artifact and none is needed --- fmix moves are
function-preserving, so the encode/compensate structure may be chewed
arbitrarily. DB-warmth is preserved by biasing realizations toward g57
shapes. One interaction to note: fmix's merge/cancel machinery may
collapse an inject/flush pair it happens to bring adjacent --- function
is unaffected, and by that point structural mixing has taken over the
hiding; the polymorphic realizations exist to protect the \emph{unmixed}
artifact from the same cancellation performed adversarially.

\textbf{Contract.} Every inject is compensated before decode, so the
emitted circuit computes the identical function: the zero-slice contract
$A(x,0) = (C(x), J)$, the pinned per-round \code{GadgetLow} check, S1's
inverse protection, and the feistel/legacy paths are all untouched. Junk
$J$ differs when flushes target the opposite carrier; $J$ is
unconstrained.

\textbf{What this does not fix} (tracked elsewhere): gate-order/positional
legibility (commuting shuffle $+$ mixing), the static gate-statistics
front (instrument still to be built), and S1's off-slice strength.

\section{Verification and measurement plan}

Unit and property tests:
\begin{enumerate}
\item \emph{Inject/flush identity}: for every realization pair and
polarity draw, inject $+$ flush with untouched sources is a wire-level
identity (exhaustive over small wire sets, like \code{cg\_menu} tests).
\item \emph{Tower exactness}: inject, apply a registered RG3/RG4 write to
a source, extend, flush; assert wire-level identity; repeat to the depth
cap.
\item \emph{Masked gadgetize end-to-end}: the pinned zero-slice check on
masked builds at several $(k, \text{cov}, \text{depth})$ points; plus a
debug-build invariant checker that simulates random inputs and asserts
$v_i = c_{i,0} \oplus c_{i,1} \oplus \bigoplus \varphi_{i,t}$ at every
unit boundary.
\item \emph{Vocabulary compliance}: no bare X; census stats within the
CG-menu envelope.
\end{enumerate}

Measurement (the arbiter), all read with the ridge measure
(\code{plot\_hmap\_ridge}: depth / $\rho$ / perm $z$ / contrast --- never
the mean):
\begin{enumerate}
\item \textbf{$n{=}16$ pre-mix grid} over $k \in \{1,2,3\}$, cov $\in
\{0.5, 1\}$, depth $\in \{1,2\}$, policies (a)/(b), vs.\ the current
recipe. Degree-1 prediction: ridge depth collapses toward the piling-up
row ($0.35 \to \lesssim 0.1$ at $k{=}1$, noise at $k{=}3$); perm $z$
loses significance.
\item \textbf{The tower validation}, the clean isolate: at $n{=}16$,
\emph{full} degree-2 regression is tractable (32 wires, 528 quadratic
regressors). Prediction: full degree-2 \emph{restores} the diagonal on
depth-1 masks (control --- the mechanism is exactly quadratic) and
\emph{fails} on depth-2 towers (payoff --- readout is now cubic). This
pins the causal story to the encoding degree with no confounds.
\item \textbf{Production $n{=}128$}: degree-1 $+$ restricted degree-2
(\code{--deg2-wires} at the campaign's settings) pre-mix, then through
the standard two-phase fmix pipeline. Prediction: the pre-mix diagonal is
gone at both, and --- the headline claim to test --- the post-mix
$\rho{=}1$ diagonal that survived 4M moves does not re-emerge, because
there is no degree-1 snapshot readout for mixing to fail to erase.
\item \textbf{Cost curve}: emitted sizes across the grid, benchmarked
like the generation-vs-size study.
\end{enumerate}

\section{Implementation sketch}

In \code{src/replace/gadgets.rs}, on the \code{gadgetize\_cnot} path:

\begin{itemize}
\item \code{MaskLedger}: \code{by\_value: Vec<Vec<MaskTerm>>},
\code{by\_source: HashMap<u16, Vec<(value, term)>>}; \code{MaskTerm}
holds the pending polynomial (list of monomials: support wires $+$
polarities), the tower depth, and nothing else --- realizations are drawn
at emission time.
\item Emission loop: wrap each unit with
\code{ledger.before\_unit(net\_writes)} (flush/extend interception) and,
per gap, \code{ledger.top\_up(census, rng)} toward
\code{--mask-cov}/\code{--mask-k}. CG reads consult
\code{--mask-read-policy}. \code{ledger.flush\_all()} before the decode
loop, with the taper handled by \code{top\_up} refusing re-injection in
the last \code{--mask-taper} gaps.
\item Net-writes per unit are static per unit kind (CG: target carrier;
legacy RG1: its four carriers; RG2: two; RG3: two; RG4 events: one);
borrowed scratch is excluded by construction.
\item The RG draw distribution gains RG4 events; alternatively RG4
maintenance runs as its own per-gap step, keeping \code{--rg-frequency}
semantics untouched --- preferred, since coverage is a target, not a
rate.
\item Flag surface on the \code{sss --cnot} driver; meters into the
existing report line.
\end{itemize}

Nothing changes in the decode bookends, S1, the CG menu, the RG networks,
or \code{commuting\_shuffle}.

\section*{Summary}

The diagonal survives mixing because the encoding hands every snapshot to
the affine adversary. The deferred-mask encoding withdraws that gift with
the cheapest possible mechanism the vocabulary admits: the RG3 move the
body already performs, split in time, tracked by a ledger, compensated
polymorphically, and stackable into towers when degree 2 is not enough.
It costs body size (measured, knob-controlled), keeps every contract and
verification of the new-gadgetize recipe, reuses the phase-A
census/targeting machinery, and comes with a falsifiable prediction
schedule --- including one experiment ($n{=}16$ full degree-2,
depth 1 vs.\ 2) that isolates the mechanism exactly.

\end{document}
